% -*- coding: utf-8 -*-
% Aula: Busca em profundidade (DFS)
% Baseada em: Paulo Feofiloff, "Algoritmos para Grafos (em linguagem C)",
% IME-USP, https://www.ime.usp.br/~pf/algoritmos_para_grafos/
\documentclass[aspectratio=169]{beamer}

\usetheme{UFS}

\usepackage[utf8]{inputenc}
\usepackage[T1]{fontenc}
\usepackage[brazil]{babel}
\usepackage{amsmath,amssymb,amsthm}
\usepackage{graphicx}
\usepackage{booktabs}
\usepackage{listings}
\usepackage{xcolor}
\usepackage{tikz}
\usetikzlibrary{arrows.meta, shapes, positioning, calc, fit, matrix,
  decorations.pathreplacing, backgrounds, chains, shapes.geometric}

\definecolor{codegreen}{rgb}{0,0.6,0}
\definecolor{codegray}{rgb}{0.5,0.5,0.5}
\definecolor{codepurple}{rgb}{0.58,0,0.82}
\definecolor{backcolour}{rgb}{0.95,0.95,0.92}

\lstdefinestyle{mystyle}{
    language=C,
    backgroundcolor=\color{backcolour},
    commentstyle=\color{codegreen},
    keywordstyle=\color{magenta},
    numberstyle=\tiny\color{codegray},
    stringstyle=\color{codepurple},
    breakatwhitespace=false,
    breaklines=true,
    captionpos=b,
    keepspaces=true,
    numbers=none,
    showspaces=false,
    showstringspaces=false,
    showtabs=false,
    tabsize=2,
    frame=single,
    rulecolor=\color{gray},
    extendedchars=true,
    basicstyle=\ttfamily\scriptsize,
    literate={ê}{{\^e}}1 {à}{{\`a}}1 {ú}{{\'u}}1 {õ}{{\~o}}1 {ó}{{\'o}}1
             {í}{{\'i}}1 {á}{{\'a}}1 {ã}{{\~a}}1 {é}{{\'e}}1 {ç}{{\c{c}}}1
             {â}{{\^a}}1 {ô}{{\^o}}1 {ü}{{\"u}}1 {Á}{{\'A}}1 {É}{{\'E}}1
}
\lstset{style=mystyle}

% ---- estilos TikZ reutilizados nos diagramas de grafos ----
\tikzset{
  vtx/.style={circle, draw=blue!70, fill=blue!10, minimum size=5.5mm,
              inner sep=0pt, font=\scriptsize},
  vtxhl/.style={circle, draw=orange!80!black, fill=orange!25, minimum size=5.5mm,
                inner sep=0pt, font=\scriptsize},
  vtxdone/.style={circle, draw=green!50!black, fill=green!15, minimum size=5.5mm,
                  inner sep=0pt, font=\scriptsize},
  arc/.style={-{Stealth[length=2mm]}, thick, blue!60!black},
  arcgray/.style={-{Stealth[length=2mm]}, thick, gray!50},
  arctree/.style={-{Stealth[length=2mm]}, very thick, orange!80!black},
  edg/.style={thick, blue!60!black},
  cell/.style={draw=gray!60, minimum size=4.2mm, inner sep=0pt, font=\tiny},
  lnode/.style={draw=blue!60, fill=blue!8, minimum height=4.5mm,
                minimum width=4.5mm, inner sep=1pt, font=\tiny},
  stk/.style={draw=gray!70, fill=gray!10, minimum width=6mm,
              minimum height=4mm, inner sep=1pt, font=\tiny},
  box/.style={rectangle, draw=blue!60, fill=blue!10, rounded corners,
              minimum width=4.6cm, minimum height=0.55cm, font=\scriptsize}
}

% ---- ponteiro de leitura no rodape do frame ----
\newcommand{\fonte}[1]{%
  \vfill
  {\tiny\color{gray!55!black}\textbf{Para aprofundar:} #1\par}%
}

\title{Busca em profundidade}
\subtitle{Algoritmos e Estruturas de Dados 2}
\author[André Y. Kashiwabara <andreyoshiaki@dcomp.ufs.br>]{Prof.\ Dr.\ André Yoshiaki Kashiwabara}
\institute{Departamento de Computação --- Universidade Federal de Sergipe}
\date{\today}

\begin{document}

\begin{frame}[plain,noframenumbering]
  \titlepage
\end{frame}

% =====================================================================
\section{Do modelo ao algoritmo}

\begin{frame}{Onde paramos}
  \begin{columns}[T]
    \begin{column}{0.48\textwidth}
      \textbf{Já temos o modelo}

      \medskip
      {\small
      \begin{itemize}\scriptsize
        \item grafo = vértices $+$ arcos;
        \item matriz de adjacências: espaço $V^2$, adjacência em tempo constante;
        \item listas de adjacência: espaço $V+A$, percurso de vizinhos
              proporcional ao grau.
      \end{itemize}
      Sabemos \emph{guardar} um grafo.}
    \end{column}
    \begin{column}{0.48\textwidth}
      \textbf{Falta o algoritmo}

      \medskip
      {\small
      Nenhuma das duas representações responde, sozinha, a perguntas como:
      \begin{itemize}\scriptsize
        \item existe caminho de $v$ até $w$?
        \item quais vértices são alcançáveis a partir de $s$?
        \item o grafo tem ciclo? em quantos pedaços ele se divide?
      \end{itemize}
      Todas exigem \emph{percorrer} o grafo.}
    \end{column}
  \end{columns}
  \fonte{representações de grafos em \cite{feofiloff2020algoritmos}, cap.\
  \emph{Estruturas de dados para grafos}, e \cite{cormen2009introduction}, §22.1.}
\end{frame}

\begin{frame}{Caminhos e alcançabilidade}
  \begin{block}{Definições}
    {\small Um \textbf{caminho} de $v$ a $w$ é uma sequência de vértices
    $v = u_0, u_1, \ldots, u_k = w$ em que $u_{i}$\texttt{-}$u_{i+1}$ é arco do
    grafo para todo $i$. Dizemos que $w$ é \textbf{alcançável} a partir de $v$
    se existe tal caminho.}
  \end{block}
  \begin{exampleblock}{Por que isso é difícil por inspeção}
    {\small O número de sequências de vértices cresce exponencialmente com $V$.
    Precisamos de um percurso que visite cada vértice \emph{no máximo uma vez} ---
    e que, mesmo assim, não deixe nenhum vértice alcançável de fora.}
  \end{exampleblock}
  \fonte{caminhos, passeios e conexidade em \cite{feofiloff2011teoriagrafos}
  (em português) e \cite{cormen2009introduction}, apêndice~B.4.}
\end{frame}

\begin{frame}{O grafo desta aula}
  {\small O mesmo da aula anterior: vértices $0..5$ e o conjunto de arcos}
  \begin{center}\ttfamily\small 0-1\quad 0-5\quad 1-0\quad 1-5\quad 2-4\quad 3-1\quad 5-3\end{center}
  \begin{center}
  \begin{tikzpicture}[scale=0.9]
    \node[vtx] (v0) at (0,1.3)  {0};
    \node[vtx] (v1) at (2.0,1.3){1};
    \node[vtx] (v2) at (4.0,1.3){2};
    \node[vtx] (v3) at (0,0)    {3};
    \node[vtx] (v5) at (2.0,0)  {5};
    \node[vtx] (v4) at (4.0,0)  {4};
    \draw[arc] (v0) to[bend left=18]  (v1);
    \draw[arc] (v1) to[bend left=18]  (v0);
    \draw[arc] (v0) -- (v5);
    \draw[arc] (v1) -- (v5);
    \draw[arc] (v2) -- (v4);
    \draw[arc] (v3) -- (v1);
    \draw[arc] (v5) -- (v3);
  \end{tikzpicture}
  \end{center}
  {\scriptsize \textbf{Convenção da aula:} os vizinhos de um vértice são sempre
  examinados em ordem crescente. Isso torna o percurso \emph{determinístico} e
  as simulações reproduzíveis.}
  \fonte{grafo-exemplo e convenções de código de \cite{feofiloff2020algoritmos}.}
\end{frame}

% =====================================================================
\section{A ideia da busca em profundidade}

\begin{frame}{A regra do labirinto}
  \begin{block}{Busca em profundidade, em uma frase}
    {\small Ao chegar a um vértice ainda não visitado, \textbf{marque-o} e siga
    imediatamente por um arco que leve a um vértice ainda não marcado; quando não
    houver nenhum, \textbf{volte} ao vértice de onde você veio e continue de lá.}
  \end{block}
  \medskip
  {\small
  \begin{itemize}\small
    \item ``Profundidade'' é a prioridade: avançar sempre que possível, recuar só
          quando não há saída nova.
    \item A marcação é o que impede voltar a andar em círculos --- o fio que
          Teseu desenrola no labirinto.
    \item A recursão da linguagem C já implementa o ``volte ao vértice de onde
          você veio'': é a pilha de chamadas.
  \end{itemize}}
  \fonte{a ideia da busca em profundidade em \cite{feofiloff2020algoritmos},
  cap.\ \emph{Busca em profundidade}; \cite{sedgewick2011algorithms}, §4.1;
  \cite{levitin2012introduction}, §3.5.}
\end{frame}

\begin{frame}{Por que a busca termina --- e por que não perde nada}
  \begin{block}{Terminação}
    {\small Um vértice é marcado no máximo uma vez, e só examinamos os arcos que
    saem de um vértice no momento em que ele é marcado. Logo cada arco é examinado
    no máximo uma vez: o processo é finito.}
  \end{block}
  \begin{block}{Corretude (alcançabilidade)}
    {\small Ao final da busca iniciada em $s$, um vértice está marcado
    \emph{se e somente se} é alcançável a partir de $s$.}
  \end{block}
  {\scriptsize \emph{Esboço}: todo vértice marcado foi atingido por uma sequência
  de arcos a partir de $s$, logo é alcançável. Reciprocamente, se algum vértice
  alcançável ficasse desmarcado, haveria no caminho de $s$ até ele um primeiro
  arco $v$\texttt{-}$w$ com $v$ marcado e $w$ desmarcado --- impossível, pois ao
  marcar $v$ examinamos todos os seus vizinhos.}
  \fonte{prova formal (teorema do caminho branco) em
  \cite{cormen2009introduction}, §22.3; a DFS como ferramenta de tempo linear
  nasce em \cite{tarjan1972depth}.}
\end{frame}

\begin{frame}[fragile]{DFS recursiva sobre a matriz de adjacências}
\begin{lstlisting}
static int cnt;
static int pre[maxV];   /* pre[v] = ordem de visita de v; -1 = nao visitado */

/* Visita todos os vertices ainda nao visitados que sao alcancaveis a
   partir de v. Supoe que G e representado por matriz de adjacencias. */
static void dfsR( Graph G, vertex v) {
   pre[v] = cnt++;                       /* marca v */
   for (vertex w = 0; w < G->V; ++w)     /* examina os vizinhos de v */
      if (G->adj[v][w] != 0)
         if (pre[w] == -1)               /* w ainda nao visitado */
            dfsR( G, w);                 /* avanca */
}                                        /* fim: volta ao chamador */
\end{lstlisting}
  {\scriptsize O vetor \texttt{pre[]} faz dois papéis ao mesmo tempo: diz
  \emph{se} o vértice já foi visitado e \emph{quando} isso aconteceu.}
  \fonte{código-fonte e variantes em \cite{feofiloff2020algoritmos};
  implementação em C em \cite{sedgewick2002algorithms}, cap.~18.}
\end{frame}

\begin{frame}[fragile]{DFS recursiva sobre listas de adjacência}
\begin{lstlisting}
static void dfsR( Graph G, vertex v) {
   pre[v] = cnt++;
   for (link a = G->adj[v]; a != NULL; a = a->next) {
      vertex w = a->w;
      if (pre[w] == -1)
         dfsR( G, w);
   }
}
\end{lstlisting}
  \begin{exampleblock}{A diferença que importa}
    {\scriptsize O corpo do algoritmo é o mesmo; muda apenas \emph{como} se enumeram
    os vizinhos: na matriz, sempre $V$ testes; nas listas, tantos passos quanto o
    grau de saída. Como \texttt{GRAPHinsertArc()} insere no \emph{início} da lista,
    os vizinhos saem na ordem inversa à de inserção e a busca a partir de 0 visita
    \texttt{0, 5, 3, 1}.}
  \end{exampleblock}
  \fonte{\cite{feofiloff2020algoritmos}; \cite{sedgewick2011algorithms}, §4.1.}
\end{frame}

% =====================================================================
\section{Simulando a busca}

\begin{frame}{Passo a passo a partir de 0 --- primeira metade}
  \begin{center}
  \begin{tikzpicture}[scale=0.74, every node/.style={transform shape}]
    \foreach \dx/\tit/\marcados in {0/{\texttt{dfsR(0)}}/{0}, 5.0/{\texttt{dfsR(1)}}/{0,1}, 10.0/{\texttt{dfsR(5)}}/{0,1,5}} {
      \begin{scope}[xshift=\dx cm]
        \draw[gray!30, rounded corners] (-0.7,-0.7) rectangle (4.3,2.0);
        \node at (1.8,1.65) {\small \tit};
        \node[vtx] (a0) at (0,1.0)  {0};
        \node[vtx] (a1) at (1.8,1.0){1};
        \node[vtx] (a2) at (3.6,1.0){2};
        \node[vtx] (a3) at (0,0)    {3};
        \node[vtx] (a5) at (1.8,0)  {5};
        \node[vtx] (a4) at (3.6,0)  {4};
        \draw[arcgray] (a0) to[bend left=18] (a1);
        \draw[arcgray] (a1) to[bend left=18] (a0);
        \draw[arcgray] (a0) -- (a5);
        \draw[arcgray] (a1) -- (a5);
        \draw[arcgray] (a2) -- (a4);
        \draw[arcgray] (a3) -- (a1);
        \draw[arcgray] (a5) -- (a3);
        \foreach \m in \marcados { \node[vtxhl] at (a\m) {\m}; }
      \end{scope}
    }
    \draw[arctree] ([xshift=5.0cm]0,1.0) -- ([xshift=5.0cm]1.8,1.0);
    \draw[arctree] ([xshift=10.0cm]0,1.0) -- ([xshift=10.0cm]1.8,1.0);
    \draw[arctree] ([xshift=10.0cm]1.8,1.0) -- ([xshift=10.0cm]1.8,0);
  \end{tikzpicture}
  \end{center}
  {\scriptsize Em \texttt{dfsR(1)} o primeiro vizinho é 0, já marcado: nada
  acontece. O segundo é 5, que ainda não foi visitado --- a busca avança.
  Laranja $=$ marcado; laranja grosso $=$ arco por onde a busca avançou.}
  \fonte{simulações passo a passo equivalentes em
  \cite{feofiloff2020algoritmos} e \cite{sedgewick2002algorithms}, cap.~18.}
\end{frame}

\begin{frame}{Passo a passo a partir de 0 --- segunda metade}
  \begin{center}
  \begin{tikzpicture}[scale=0.74, every node/.style={transform shape}]
    \foreach \dx/\tit/\marcados in {0/{\texttt{dfsR(3)}}/{0,1,5,3}, 5.0/{volta a \texttt{dfsR(5)}}/{0,1,5,3}, 10.0/{fim de \texttt{dfsR(0)}}/{0,1,5,3}} {
      \begin{scope}[xshift=\dx cm]
        \draw[gray!30, rounded corners] (-0.7,-0.7) rectangle (4.3,2.0);
        \node at (1.8,1.65) {\small \tit};
        \node[vtx] (b0) at (0,1.0)  {0};
        \node[vtx] (b1) at (1.8,1.0){1};
        \node[vtx] (b2) at (3.6,1.0){2};
        \node[vtx] (b3) at (0,0)    {3};
        \node[vtx] (b5) at (1.8,0)  {5};
        \node[vtx] (b4) at (3.6,0)  {4};
        \draw[arcgray] (b0) to[bend left=18] (b1);
        \draw[arcgray] (b1) to[bend left=18] (b0);
        \draw[arcgray] (b0) -- (b5);
        \draw[arcgray] (b1) -- (b5);
        \draw[arcgray] (b2) -- (b4);
        \draw[arcgray] (b3) -- (b1);
        \draw[arcgray] (b5) -- (b3);
        \foreach \m in \marcados { \node[vtxhl] at (b\m) {\m}; }
        \draw[arctree] (0,1.0) -- (1.8,1.0);
        \draw[arctree] (1.8,1.0) -- (1.8,0);
        \draw[arctree] (1.8,0) -- (0,0);
      \end{scope}
    }
  \end{tikzpicture}
  \end{center}
  {\scriptsize O único vizinho de 3 é 1, já marcado: a busca recua até 0, cujo
  segundo vizinho (5) também já está marcado. \textbf{Resultado:} os vértices
  alcançáveis a partir de 0 são $\{0,1,5,3\}$; \textbf{2 e 4 não são alcançáveis}.}
  \fonte{alcançabilidade e busca a partir de uma raiz em
  \cite{feofiloff2020algoritmos}; \cite{cormen2009introduction}, §22.3.}
\end{frame}

\begin{frame}{A ordem de visita: o vetor \texttt{pre[]}}
  \begin{center}
  {\small
  \begin{tabular}{lcccccc}
    \toprule
    vértice $v$   & 0 & 1 & 2 & 3 & 4 & 5 \\
    \midrule
    \texttt{pre[v]} & 0 & 1 & $-1$ & 3 & $-1$ & 2 \\
    \bottomrule
  \end{tabular}}
  \end{center}
  \medskip
  {\small
  \begin{itemize}\small
    \item A ordem $0, 1, 5, 3$ em que os vértices receberam número é a
          \textbf{ordem de pré-ordem} (ou \emph{ordem de descoberta}).
    \item Ela é atribuída na \emph{entrada} da chamada recursiva, antes de
          examinar os vizinhos.
    \item Os vértices que permanecem com $-1$ são exatamente os
          \emph{não alcançáveis} a partir da raiz da busca.
  \end{itemize}}
  \fonte{pré-ordem e carimbos de tempo (\emph{timestamps} $d[v]$ e $f[v]$) em
  \cite{cormen2009introduction}, §22.3; numeração \texttt{pre[]} em
  \cite{feofiloff2020algoritmos}.}
\end{frame}

% =====================================================================
\section{Varredura completa e floresta DFS}

\begin{frame}[fragile]{Uma busca não basta: a varredura}
\begin{lstlisting}
/* Visita TODOS os vertices do grafo, em uma ou mais buscas. */
void GRAPHdfs( Graph G) {
   cnt = 0;
   for (vertex v = 0; v < G->V; ++v)
      pre[v] = -1;
   for (vertex v = 0; v < G->V; ++v)
      if (pre[v] == -1)      /* v ainda nao alcancado por busca anterior */
         dfsR( G, v);        /* comeca uma nova busca com raiz v */
}
\end{lstlisting}
  {\scriptsize Cada iteração do segundo laço que dispara \texttt{dfsR()} inicia
  uma \textbf{nova busca}, com uma nova \textbf{raiz}. No nosso grafo há duas:
  o vértice 0 e o vértice 2.}
  \fonte{a varredura completa (\texttt{DFS}) em \cite{cormen2009introduction},
  §22.3, e em \cite{feofiloff2020algoritmos}.}
\end{frame}

\begin{frame}{A floresta DFS}
  \begin{columns}[T]
    \begin{column}{0.52\textwidth}
      \begin{center}
      \begin{tikzpicture}[scale=0.78, every node/.style={transform shape}]
        \node[vtxhl] (t0) at (0,2.1)   {0};
        \node[vtxhl] (t1) at (0,1.2)   {1};
        \node[vtxhl] (t5) at (0,0.3)   {5};
        \node[vtxhl] (t3) at (0,-0.6)  {3};
        \draw[arctree] (t0) -- (t1);
        \draw[arctree] (t1) -- (t5);
        \draw[arctree] (t5) -- (t3);
        \node[vtxhl] (t2) at (2.4,2.1) {2};
        \node[vtxhl] (t4) at (2.4,1.2) {4};
        \draw[arctree] (t2) -- (t4);
        \node[font=\tiny] at (0.9,2.1)  {pre $=0$};
        \node[font=\tiny] at (0.9,1.2)  {1};
        \node[font=\tiny] at (0.9,0.3)  {2};
        \node[font=\tiny] at (0.9,-0.6) {3};
        \node[font=\tiny] at (3.3,2.1)  {4};
        \node[font=\tiny] at (3.3,1.2)  {5};
      \end{tikzpicture}
      \end{center}
    \end{column}
    \begin{column}{0.44\textwidth}
      {\small
      Os arcos por onde a busca \emph{avançou} formam uma floresta:
      \begin{itemize}\scriptsize
        \item cada árvore tem por raiz um vértice em que o laço de
              \texttt{GRAPHdfs()} disparou uma busca;
        \item os demais arcos do grafo (como 1\texttt{-}0 e 3\texttt{-}1)
              \emph{não} entram na floresta --- eles levam a vértices já marcados;
        \item a floresta tem $V$ vértices e $V - (\text{número de raízes})$ arcos.
      \end{itemize}
      \textbf{Na próxima aula} classificaremos justamente os arcos que ficaram
      de fora.}
    \end{column}
  \end{columns}
  \fonte{floresta DFS e classificação de arcos em
  \cite{cormen2009introduction}, §22.3; \cite{sedgewick2002algorithms},
  cap.~18; origem em \cite{tarjan1972depth}.}
\end{frame}

% =====================================================================
\section{A versão iterativa}

\begin{frame}[fragile]{DFS com pilha explícita}
\begin{lstlisting}
void GRAPHdfsPilha( Graph G, vertex s) {
   for (vertex v = 0; v < G->V; ++v) pre[v] = -1;   cnt = 0;
   STACKinit( G->V * G->V);            /* pilha de vertices */
   STACKpush( s);
   while (!STACKempty()) {
      vertex v = STACKpop();
      if (pre[v] != -1) continue;      /* ja visitado por outro caminho */
      pre[v] = cnt++;
      for (vertex w = 0; w < G->V; ++w)
         if (G->adj[v][w] != 0 && pre[w] == -1)
            STACKpush( w);
   }
}
\end{lstlisting}
  {\scriptsize A pilha substitui a pilha de chamadas. O teste \texttt{pre[v] != -1}
  após o \texttt{pop} é necessário: um vértice pode ser empilhado várias vezes.}
  \fonte{pilha explícita: \cite{sedgewick2002algorithms}, cap.~18; \cite{ziviani2011projeto}, §7.3.}
\end{frame}

\begin{frame}{Mesma busca, outra ordem}
  \begin{columns}[T]
    \begin{column}{0.46\textwidth}
      {\small \textbf{Recursiva, a partir de 0}}

      \medskip
      {\scriptsize Ordem de visita: \texttt{0, 1, 5, 3}

      \medskip
      Os vizinhos são examinados em ordem crescente e cada avanço acontece
      \emph{antes} de o vizinho seguinte ser considerado.}
    \end{column}
    \begin{column}{0.50\textwidth}
      {\small \textbf{Iterativa, a partir de 0}}

      \medskip
      {\scriptsize Ordem de visita: \texttt{0, 5, 3, 1}

      \medskip
      Empilhamos 1 e depois 5; o \texttt{pop} devolve \emph{o último}, isto é,
      5. O 1 fica esperando no fundo da pilha.}
    \end{column}
  \end{columns}
  \medskip
  \begin{exampleblock}{O que \emph{não} muda}
    {\small O conjunto de vértices visitados é o mesmo: $\{0,1,5,3\}$. A busca em
    profundidade não define uma ordem única --- ela depende da representação e da
    implementação. O que ela garante é \emph{quais} vértices são alcançados.}
  \end{exampleblock}
  \fonte{comparação entre as ordens recursiva e com pilha em
  \cite{sedgewick2002algorithms}, cap.~18, e \cite{sedgewick2011algorithms}, §4.1.}
\end{frame}

% =====================================================================
\section{Custo}

\begin{frame}{Quanto custa a busca}
  \begin{center}
  {\small
  \begin{tabular}{lll}
    \toprule
    & \textbf{Matriz de adjacências} & \textbf{Listas de adjacência} \\
    \midrule
    Marcar todos os vértices        & $\Theta(V)$   & $\Theta(V)$ \\
    Examinar os vizinhos de $v$     & $\Theta(V)$   & $\Theta(\text{grau de saída de } v)$ \\
    \textbf{Total da varredura}     & $\Theta(V^2)$ & $\Theta(V + A)$ \\
    \bottomrule
  \end{tabular}}
  \end{center}
  \medskip
  {\small
  \begin{itemize}\small
    \item Com listas o total sai da soma $\sum_v \text{grau de saída}(v) = A$:
          cada arco é examinado exatamente uma vez.
    \item Em grafos esparsos ($A \ll V^2$) a diferença é enorme --- e é por isso
          que quase todo algoritmo de grafos deste curso pressupõe listas.
    \item O espaço extra é $\Theta(V)$ nos dois casos: os vetores
          \texttt{pre[]} e \texttt{pai[]}, mais a pilha de recursão.
  \end{itemize}}
  \fonte{análise $\Theta(V+A)$ por agregação em \cite{cormen2009introduction},
  §22.3; discussão de esparsidade em \cite{levitin2012introduction}, §3.5, e
  \cite{feofiloff2020algoritmos}.}
\end{frame}

% =====================================================================
\section{Aplicação: recuperar um caminho}

\begin{frame}[fragile]{Guardando de onde viemos}
\begin{lstlisting}
static int pai[maxV];   /* pai[w] = vertice de onde a busca chegou a w */

static void dfsR( Graph G, vertex v) {
   pre[v] = cnt++;
   for (vertex w = 0; w < G->V; ++w)
      if (G->adj[v][w] != 0 && pre[w] == -1) {
         pai[w] = v;                  /* registra o arco de arvore v-w */
         dfsR( G, w);
      }
}
\end{lstlisting}
  {\scriptsize O vetor \texttt{pai[]} é a floresta DFS armazenada em um vetor:
  \texttt{pai[v] == v} identifica as raízes. Custa $\Theta(V)$ de espaço e nada
  de tempo adicional.}
  \fonte{vetor de pais em \cite{feofiloff2020algoritmos}; equivalentes
  $\pi[v]$ em \cite{cormen2009introduction}, §22.3, e \texttt{edgeTo[]} em
  \cite{sedgewick2011algorithms}, §4.1.}
\end{frame}

\begin{frame}{De \texttt{pai[]} ao caminho}
  \begin{columns}[T]
    \begin{column}{0.46\textwidth}
      {\small \textbf{O vetor, após a busca em 0}}
      \begin{center}\scriptsize
      \begin{tabular}{lcccc}
        \toprule
        $v$        & 1 & 5 & 3 \\
        \texttt{pai[v]} & 0 & 1 & 5 \\
        \bottomrule
      \end{tabular}
      \end{center}
      \medskip
      {\scriptsize Para achar o caminho de 0 a 3, siga \texttt{pai[]} a partir de
      3 até a raiz e inverta: $3 \to 5 \to 1 \to 0$, logo o caminho é
      \texttt{0, 1, 5, 3}.}
    \end{column}
    \begin{column}{0.50\textwidth}
      \begin{block}{Atenção}
        {\scriptsize O caminho obtido é \emph{um} caminho, não necessariamente o
        mais curto: a busca em profundidade encontrou \texttt{0,1,5,3}, mas
        \texttt{0,5,3} também é caminho e é mais curto. Caminhos mínimos em
        número de arcos são tarefa da busca em \textbf{largura}.}
      \end{block}
    \end{column}
  \end{columns}
  \fonte{recuperação de caminhos (\texttt{pathTo()}) em
  \cite{sedgewick2011algorithms}, §4.1; caminhos mínimos por busca em largura
  em \cite{cormen2009introduction}, §22.2.}
\end{frame}

% =====================================================================
\section{Fechamento}

\begin{frame}{Para praticar}
  \begin{block}{Exercícios}
    {\small
    \begin{enumerate}\small
      \item Escreva \texttt{GRAPHreach(G, s)} que devolve o número de vértices
            alcançáveis a partir de $s$.
      \item Adapte \texttt{dfsR()} para listas e compare a ordem de visita com a
            da matriz.
      \item Implemente a versão iterativa com pilha e verifique a ordem
            \texttt{0, 5, 3, 1}.
      \item Escreva \texttt{GRAPHpath(G, s, t)} que imprime um caminho de $s$ a
            $t$, ou avisa que não existe.
      \item Quantas raízes tem a floresta DFS de um grafo sem arcos? E de um
            grafo em que todo vértice alcança todos os outros?
    \end{enumerate}}
  \end{block}
  {\scriptsize O tutorial em Jupyter traz esqueletos em C para os itens 1, 3 e 4.}
  \fonte{mais exercícios em \cite{feofiloff2020algoritmos} e \cite{ziviani2011projeto}, §7.3.}
\end{frame}

\begin{frame}{Recapitulando a aula}
  \begin{block}{As cinco ideias centrais}
    \begin{itemize}\small
      \item Percorrer é o que transforma a representação em resposta: caminhos,
            alcançabilidade, ciclos.
      \item A regra da DFS: marque, avance enquanto puder, recue quando não puder.
      \item \texttt{pre[]} marca \emph{e} numera; \texttt{GRAPHdfs()} varre o
            grafo todo, gerando uma \textbf{floresta}.
      \item Custo $\Theta(V+A)$ com listas, $\Theta(V^2)$ com matriz.
      \item \texttt{pai[]} guarda a floresta em um vetor e devolve um caminho ---
            não necessariamente o mais curto.
    \end{itemize}
  \end{block}
  {\small \textbf{Próxima aula:} os arcos que ficaram fora da floresta ---
  pós-ordem, classificação de arcos, detecção de ciclos e componentes.}
\end{frame}

\begin{frame}{Roteiro de leitura}
  \begin{center}
  {\scriptsize
  \begin{tabular}{ll}
    \toprule
    \textbf{Conceito} & \textbf{Onde ler} \\
    \midrule
    Representações e o grafo-exemplo    & \cite{feofiloff2020algoritmos}; \cite{cormen2009introduction}, §22.1 \\
    Caminhos, alcançabilidade           & \cite{feofiloff2011teoriagrafos}; \cite{cormen2009introduction}, ap.~B.4 \\
    A ideia da DFS                      & \cite{feofiloff2020algoritmos}; \cite{levitin2012introduction}, §3.5 \\
    Código em C (matriz e listas)       & \cite{feofiloff2020algoritmos}; \cite{sedgewick2002algorithms}, cap.~18 \\
    Corretude e pré-ordem               & \cite{cormen2009introduction}, §22.3 \\
    Varredura e floresta DFS            & \cite{cormen2009introduction}, §22.3; \cite{tarjan1972depth} \\
    Versão iterativa com pilha          & \cite{sedgewick2002algorithms}, cap.~18; \cite{ziviani2011projeto}, §7.3 \\
    Custo $\Theta(V+A)$                 & \cite{cormen2009introduction}, §22.3; \cite{levitin2012introduction}, §3.5 \\
    Caminho a partir de \texttt{pai[]}  & \cite{sedgewick2011algorithms}, §4.1 \\
    \bottomrule
  \end{tabular}}
  \end{center}
  {\scriptsize Leitura mínima da aula: \cite{feofiloff2020algoritmos}, capítulo
  \emph{Busca em profundidade}. Os demais textos servem para aprofundar cada
  ponto.}
\end{frame}

% =====================================================================
\section*{Referências}
\begin{frame}[allowframebreaks]{Referências}
  \scriptsize
  \nocite{*}
  \bibliographystyle{plain}
  \bibliography{../referencias}
\end{frame}

\end{document}
